\subsection{Метод Гриффиса и Стюарта}
\label{subsec:griffith-stewart}

\textbf{Идея:} построить линейные аппроксимации \emph{и} для функции цели,
\emph{и} для ограничений задачи, а затем ограничить длину шага, чтобы
линеаризация оставалась осмысленной.

Метод разбирается на примере.

\begin{example}\label{ex:griffith}
\[
  (\mathrm{I})
  \qquad
  \begin{cases}
    \varphi_0(x) = x_1^2 + x_2^2 - 16 x_1 - 10 x_2 \to \min , \\[2pt]
    g_1(x_1, x_2) = x_1^2 - 6 x_1 + 4 x_2 - 11 \le 0 , \\[2pt]
    g_2(x_1, x_2) = e^{\,x_1 - 3} - x_1 x_2 + 3 x_2 - 1 \le 0 , \\[2pt]
    x_1 \ge 0 , \quad x_2 \ge 0 .
  \end{cases}
\]
\end{example}

\subsubsection*{Начальный шаг, $k = 0$}

\[
  x^{(0)} = \begin{pmatrix} 4 \\ 3 \end{pmatrix} ,
  \qquad
  \varphi_0\big(x^{(0)}\big) = -69 .
\]
Линеаризуем все функции в точке $x^{(0)}$. Для целевой функции
\[
  \tilde\varphi_0(x) = -69 +
  \begin{pmatrix} 2x_1 - 16 \\ 2x_2 - 10 \end{pmatrix}\T_{\!\!\big|\,x^{(0)}}
  \begin{pmatrix} x_1 - 4 \\ x_2 - 3 \end{pmatrix}
  = -69 - 8(x_1 - 4) - 4(x_2 - 3)
  = -8x_1 - 4x_2 - 25 ,
\]
и аналогично для ограничений. Получается задача ЛП
\[
  (\mathrm{II})
  \qquad
  \begin{cases}
    \tilde\varphi_0(x) = -8x_1 - 4x_2 - 25 \to \min , \\[2pt]
    \tilde g_1(x) = 2x_1 + 4x_2 - 27 \le 0 , \\[2pt]
    \tilde g_2(x) = -0{,}28 x_1 - x_2 + 2{,}84 \le 0 , \\[2pt]
    x_1 \ge 0 , \quad x_2 \ge 0 .
  \end{cases}
\]
Её решение —
\[
  x^{(1)} = \begin{pmatrix} 13{,}5 \\ 0 \end{pmatrix} ,
  \qquad\text{тогда как}\qquad
  x^{(0)} = \begin{pmatrix} 4 \\ 3 \end{pmatrix} .
\]

Точка ушла слишком далеко: линеаризация в окрестности $x^{(0)}$ ничего не знает
про поведение $g_1, g_2$ при $x_1 \approx 13{,}5$. Поэтому накладывается
ограничение на шаг, предотвращающее выход из области:
\begin{equation}\label{eq:gs-trust}
  \big| x_j^{(1)} - x_j^{(0)} \big| \ \le\ \delta_j^{(0)} ,
  \qquad j = \overline{1,2} .
\end{equation}
Положив $\delta^{(0)} = (0{,}5;\ 0{,}5)\T$, получаем
\[
  x^{(1)} = \begin{pmatrix} 4 + 0{,}5 \\ 3 - 0{,}5 \end{pmatrix}
          = \begin{pmatrix} 4{,}5 \\ 2{,}5 \end{pmatrix} .
\]

\subsubsection*{Шаг $k = 1$}

Строим линейную аппроксимацию задачи~$(\mathrm{I})$ в окрестности $x^{(1)}$:
\[
  (\mathrm{III})
  \qquad
  \begin{cases}
    \tilde\varphi_0(x) = -7 x_1 - 5 x_2 - 26{,}5 \to \min , \\[2pt]
    3 x_1 + 4 x_2 - 31{,}40 \le 0 , \\[2pt]
    1{,}98 x_1 - 1{,}50 x_2 - 5{,}42 \le 0 , \\[2pt]
    x_1 \ge 0 , \quad x_2 \ge 0 .
  \end{cases}
\]
Решение ЛП —
\[
  x^{(2)} = \begin{pmatrix} 5{,}45 \\ 3{,}58 \end{pmatrix}
  \qquad\text{— не подходит, снова применяем \eqref{eq:gs-trust}:}
\]
\[
  \delta_j^{(1)} = \lambda\, \delta_j^{(0)} , \qquad \lambda \in (0,1) .
\]
При $\delta^{(1)} = (0{,}4;\ 0{,}4)\T$
\[
  x^{(2)} = \begin{pmatrix} 4{,}90 \\ 2{,}90 \end{pmatrix} ,
  \qquad
  \varphi_0\big(x^{(2)}\big) = -75{,}1 .
\]

\begin{idea}[title={О числах в примере}]
Свободный член линеаризованной цели ($-25$ на первом шаге, $-26{,}5$ на втором)
на решение задачи ЛП не влияет — минимум линейной функции не меняется от сдвига
на константу. Поэтому расхождения в этом слагаемом несущественны; важны только
коэффициенты при переменных, то есть компоненты градиента.

Коэффициенты в $\tilde g_1, \tilde g_2$ получаются прямой подстановкой:
\[
  \nabla g_1 = \begin{pmatrix} 2x_1 - 6 \\ 4 \end{pmatrix} ,
  \qquad
  \nabla g_2 = \begin{pmatrix} e^{\,x_1-3} - x_2 \\ -x_1 + 3 \end{pmatrix} ,
\]
например в $x^{(0)} = (4,3)$: $\nabla g_1 = (2;\,4)\T$,
$g_1(x^{(0)}) = -7$, откуда $\tilde g_1 = -7 + 2(x_1-4) + 4(x_2-3)
= 2x_1 + 4x_2 - 27$; и $\nabla g_2 = (e - 3;\, -1)\T \approx (-0{,}28;\,-1)\T$,
$g_2(x^{(0)}) \approx -1{,}28$, откуда
$\tilde g_2 \approx -0{,}28x_1 - x_2 + 2{,}84$. Числа в $(\mathrm{III})$
отличаются от пересчёта в третьем знаке — это округление, на ход метода оно не
влияет.
\end{idea}

\begin{idea}[title={Смысл ограничения \eqref{eq:gs-trust}}]
Линеаризация точна только в малой окрестности точки разложения, а минимум
линейной функции на многограннике всегда достигается в вершине — то есть
«как можно дальше». Без дополнительного ограничения метод на каждом шаге
улетал бы в далёкую вершину, где аппроксимация уже не имеет отношения к
исходной задаче (так и вышло: $x^{(1)} = (13{,}5;\,0)$ вместо разумного
смещения от $(4;\,3)$).

Неравенства~\eqref{eq:gs-trust} задают вокруг текущей точки куб со стороной
$2\delta_j^{(k)}$ — \emph{область доверия} к линейной модели. Если шаг оказался
неудачным, размеры куба уменьшают: $\delta^{(k+1)} = \lambda\delta^{(k)}$,
$\lambda \in (0,1)$ — то же самое дробление, что и при выборе шага в
градиентных методах, только дробится не длина шага вдоль фиксированного
направления, а размер всей области, в которой ищется следующая точка.
\end{idea}
